\begin{theorem}[Trace on the split elliptic normalization]\hypertarget{tag-0308}{}\label{push:src:B.3.4}\label{push:B34}
\StatementRecord{push:src:B.3.4}{proved}{0308}{\SeriesRef{push:src:B.3.2}; \cite[Theorems 8.4 and 8.6, table (8.16)]{Shioda90}; \cite[Theorem 4.1]{Shioda90}; \cite[Theorem 1.3]{Shioda90}}

Let $C$ be the genus-three component cover
in \ref{push:src:B.3.2}, and let $\mathscr E\to C$ be
the relatively minimal elliptic surface
obtained by pulling back $S\to B$.
Write $\sigma$ for the order-twelve deck
transformation, with its natural lift to
$\mathscr E$. Then
\[
 \operatorname{Tr}_{\sigma}
   \mathscr E(\mathbb C(C))
 :=\left\{\sum_{i=0}^{11}\sigma^i Q:
                 Q\in\mathscr E(\mathbb C(C))\right\}
 =12\mathbb Z P.
\]
In particular there is no section $Q$
for which $(t_Q\circ\sigma)^{12}=t_{6P}$.
For a normalization over $C$ ruled over this split elliptic
surface, the twelfth iterate of an adjacent conductor map
$t_Q\circ\sigma$ is therefore $t_{12kP}$ for some $k\in\mathbb Z$.

\end{theorem}


\begin{proof}
\begin{enumerate}[label=\arabic*.,leftmargin=*]
\item There is a convenient model of the cover
and its elliptic surface:
\[
 C:\ b^3=a^4+1,\qquad
 t=b^3,\quad y=b/a,\qquad
 \mathscr E:\ Y^2=X^3-3bX+2a^2.
\]
Indeed $y^{12}=t^4/(t-1)^3$.
The substitution from the original
Weierstrass coordinates is
$x=a^2b^4X$, $y_{\rm old}=a^3b^6Y$.
The affine discriminant of the new
equation is $1728$; it has no singular
fibre on the affine curve. The curve
$C$ has one point at infinity, with
pole orders $3,4$ for $a,b$.
Put $z=a/b$ and $c=z^3a=z^4b$ there.
Then $z$ is a uniformizer and
\[
 c^3(1-c)=z^{12},\qquad c(0)=1.
\]
In the integral Weierstrass coordinates
$X_0=z^2X$, $Y_0=z^3Y$, the equation is
\[
 Y_0^2=X_0^3-3cX_0+2c^2.
\]
Its discriminant has order twelve and
its $c_4$ is a unit. The unique singular
fibre of the minimal regular surface
is therefore $I_{12}$. The fundamental
line bundle is $\mathcal O_C(\infty)$,
so $\chi(\mathcal O_{\mathscr E})=1$.
This is a projective elliptic surface.

\item Let $\omega=e^{2\pi i/3}$ and
$\lambda=e^{2\pi i/12}$. The deck lift is
\[
 \sigma(a,b,X,Y)=(ia,\omega b,\lambda^2X,\lambda^3Y).
\]
It fixes the section $P=(0,-\sqrt2a)$.
Two further sections are
\[
 R=\left(\frac{a^2}{b-1},
                 \frac{\sqrt3a}{b-1}\right),
 \qquad
 R'=\sigma R=
   \left(\frac{a^2}{b-\omega},
                 \frac{\sqrt3\omega a}{b-\omega}\right).
\]
Direct substitution verifies their
equations. The three orbit points
$R,\sigma R,\sigma^2R$ lie on the line
\[
 Y=\frac{\sqrt3b}{a}X-\sqrt3a.
\]
Thus $R+\sigma R+\sigma^2R=O$ by the
elliptic chord law, and $\sigma^3R=R$.


\item  The only pole of $R$ on the affine
curve is at $(a,b)=(0,1)$. There $a$ is
a uniformizer and $b-1=a^4/3+O(a^8)$,
so $X=3a^{-2}+O(a^2)$ and
$Y=3\sqrt3a^{-3}+O(a)$.
Consequently $R\cdot O=1$.
The same computation at $(0,\omega)$
gives $R'\cdot O=1$, while $P\cdot O=0$.
At the other two points with $a=0$,
$P$ and $R$ meet simply: the fibre is
smooth at $(X,Y)=(0,0)$, and the
difference of their $Y$ coordinates
has a nonzero linear term in $a$.
Thus $P\cdot R=2$, and likewise
$P\cdot R'=2$. The sections $R,R'$
meet simply at $(a,b)=(0,\omega^2)$.
They cannot meet elsewhere on the
affine curve, since their $X$
coordinates are distinct for $a\ne0$.

\item To check infinity, for the section
$\sigma^jR$ put $k=\omega^j$. Its
integral coordinates are
\[
 X_0=\frac{c^2}{c-kz^4},\qquad
 Y_0=\frac{\sqrt3kc z^4}{c-kz^4}.
\]
At the node use analytic coordinates
\[
 U=Y_0-(X_0-1)\sqrt{X_0+2},\qquad
 V=Y_0+(X_0-1)\sqrt{X_0+2}.
\]
They satisfy
\[
 UV=(1-c)\{3X_0-2(c+1)\},
\]
whose right side is a unit times $z^{12}$.
The section expansions are
\[
 U=-\frac{\sqrt3k^2}{6}z^8+O(z^{12}),
 \qquad V=2\sqrt3kz^4+O(z^8).
\]
Hence $R,R'$ both meet component $4$
of the resolved $I_{12}$ cycle, with
distinct coordinates on that component.
Reversing the cycle numbering replaces
both indices by $8$ and has no effect
on the argument. The section $P$ meets
the identity component, at the smooth
point $(X_0,Y_0)=(0,-\sqrt2)$.
Thus no further intersections occur at
infinity, and $R\cdot R'=1$.

\item The height pairing formula and its
positivity are those of
\cite[Theorems~8.4 and~8.6, table~(8.16), pp.~227--229]{Shioda90}.
For an $I_{12}$ fibre its correction is
$\operatorname{contr}(i,j)=\min(i,j)-ij/12$.
The preceding intersections give, in
the ordered basis $P,R,R'$,
\[
 G=\begin{pmatrix}
 2&0&0\\
 0&4/3&-2/3\\
 0&-2/3&4/3
 \end{pmatrix},\qquad \det G=8/3.
\]
In particular these are three independent
sections. The Leray sequence and
$R^1\pi_*\mathcal O_{\mathscr E}
=\mathcal O_C(-\infty)$ show that its
irregularity equals $g(C)=3$; the
right-hand sheaf has no global section.
This also agrees with the Picard-variety
statement in \cite[Theorem~4.1, pp.~219--220]{Shioda90}.
The equalities $\chi=1$ and
$e(\mathscr E)=12$ give $p_g=3$,
$b_2=22$ and $h^{1,1}=16$.
The zero section, a fibre and the eleven
nonidentity $I_{12}$ components generate
a rank-thirteen trivial lattice.
The identification of sections with
the N\'eron--Severi quotient
\cite[Theorem~1.3, p.~214]{Shioda90}, together
with $\rho\le h^{1,1}$, now shows that
the Mordell--Weil rank is exactly three
and $\rho=16$.

\item There is no nonzero torsion section.
Indeed its height would be zero, so
for some $0\le j<12$ the height formula
would give
\[
 0=2+2(Q\cdot O)-\frac{j(12-j)}{12}.
\]
The correction is at most three and
$Q\cdot O$ is a nonnegative integer.
This forces $Q\cdot O=0$ and
$j(12-j)=24$, which has no integral
solution.

\item Let $L=\mathbb ZP+\mathbb ZR+\mathbb ZR'$.
The integral N\'eron--Severi sublattice
generated by the trivial lattice and
these three sections has determinant
$-32$. This follows either by eliminating
the trivial lattice with the displayed
height matrix, or from the integral
intersection matrix specified by the following entries.
All its entries have just been computed:
sections have square $-1$, their mutual
intersections are $2,2,1$, and $R,R'$
meet precisely component $4$.
If $d=[\mathscr E(\mathbb C(C)):L]$,
the section--N\'eron--Severi quotient
identification gives the same index
for this sublattice in
$\operatorname{NS}(\mathscr E)$.
Its integral determinant implies
$d^2\mid32$, hence $d\in\{1,2,4\}$.
In particular $4Q\in L$ for every section.

\item In the rational basis $P,R,R'$ the
deck action and its trace are
\[
 M=\begin{pmatrix}1&0&0\\0&0&-1\\0&1&-1\end{pmatrix},
 \qquad \sum_{i=0}^{11}M^i=
       \operatorname{diag}(12,0,0).
\]
Suppose first that the trace of $Q$ is
$6P$. By the index bound we may write
\[
 Q=\frac12P+\frac14(uR+vR'),\qquad u,v\in\mathbb Z.
\]
Let $j\in\mathbb Z/12$ be its component
index at infinity. Since both $R,R'$
meet component four, the height formula
modulo integers gives
\[
 \frac{2u-v}{6}\equiv\frac j3\pmod{\mathbb Z},
 \qquad
 \frac{-u+2v}{6}\equiv\frac j3\pmod{\mathbb Z}.
\]
Equivalently $2u-v-2j$ and $-u+2v-2j$
are divisible by six. Reducing modulo
two shows that both $u,v$ are even.
Write $u=2u'$, $v=2v'$. The component
map is a group homomorphism, so
\[
 2j\equiv4(u'+v')\pmod{12}.
\]
Therefore $j$ is even. The geometric
height formula then lies in
$\tfrac13\mathbb Z$. In contrast the
height from $G$ is
\[
 \langle Q,Q\rangle
 =\frac12+\frac{u'^2-u'v'+v'^2}{3},
\]
which lies in the disjoint coset
$\tfrac12+\tfrac13\mathbb Z$.
This is a contradiction.

\item The same argument determines the full
trace image. If $Q=\alpha P+\beta R+\gamma R'$
in the rational span, pairing with the
narrow section $P$ gives
$2\alpha=\langle Q,P\rangle\in\mathbb Z$.
Every odd half-integral value of $\alpha$
leads to the identical congruences and
height contradiction, with $1/2$
replaced by an odd square divided by
two. Thus $\alpha\in\mathbb Z$ and
the trace is in $12\mathbb ZP$.
The section $P$ realizes the generator.

Finally the identity
$(t_Q\circ\sigma)^{12}=t_{\operatorname{Tr}_\sigma Q}$
proves the claimed obstruction.
\end{enumerate}
\end{proof}
